% =====================================================================
\section{Плата как клиент: HTTP, JSON, время и файлы}
% =====================================================================
\lessonintro
  {получать данные из интернета, разбирать ответ в формате JSON, узнавать точное время и хранить файлы настроек во Flash}
  {урок 5 (Wi-Fi), урок 11 (NVS и Flash), урок 13 (время для TLS)}
  {плата запрашивает данные у сервера, показывает московское время и сохраняет файл настроек}
  {доступ к Wi-Fi с выходом в интернет; библиотека ArduinoJson версии 7}

\subsection{Сначала: плата ходит в интернет, как браузер}

В уроке 5 плата была сервером: к ней приходили запросы. Теперь поменяем роли --- плата сама отправляет
запрос и ждёт ответ. Для человека сервер вернул бы красивую страницу, а для программы --- короткий текст
с данными в формате JSON: «город, температура, влажность». Дальше плата вытаскивает из текста нужные числа.
А чтобы знать, какой сейчас час, она спрашивает точное время у специальных серверов в сети.

\subsection{Путь запроса}

Перед запросом нужно превратить имя сайта в IP-адрес. Это делает DNS-сервер, адрес которого плата
получает от роутера. Затем плата устанавливает соединение с сервером и отправляет HTTP-запрос
\code{GET}. Сервер отвечает кодом состояния (200 означает «всё хорошо») и телом ответа.

\begin{figure}[H]
\centering
\begin{tikzpicture}[y=-0.8cm,font=\sffamily\scriptsize]
  \node[blk,fill=blkcpu,minimum width=2.6cm] (e) at (0,0) {ESP32-S3};
  \node[blk,fill=blkrf,minimum width=2.6cm] (d) at (6.2,0) {DNS};
  \node[blk,fill=blkrf,minimum width=2.6cm] (s) at (12.4,0) {Сервер API};
  \draw[thick,gray] (e.south) -- (0,7.6);
  \draw[thick,gray] (d.south) -- (6.2,7.6);
  \draw[thick,gray] (s.south) -- (12.4,7.6);
  \draw[arr,accent] (0,1.3) -- node[above]{1. адрес example.com?} (6.2,1.3);
  \draw[arr,okgreen] (6.2,2.3) -- node[above]{IP-адрес} (0,2.3);
  \draw[arr,accent] (0,3.5) -- node[above,pos=0.78]{2. GET /api/weather.json} (12.4,3.5);
  \draw[arr,okgreen] (12.4,4.6) -- node[above,pos=0.78]{3. 200 OK + JSON-текст} (0,4.6);
  \fill[esp!30] (-0.15,4.9) rectangle (0.15,6.4);
  \node[anchor=west,align=left] at (0.35,5.65) {4. \code{deserializeJson()}:\\ текст $\to$ переменные};
  \node[anchor=west] at (0.35,7.1) {5. \code{http.end()}: освободить соединение};
\end{tikzpicture}
\caption{Путь запроса: имя превращается в IP-адрес, затем идёт запрос, ответ и разбор}
\end{figure}

\subsection{HTTP-запрос: GET}

Клиент \code{HTTPClient} входит в ядро Arduino-ESP32, ставить библиотеку не нужно. Внизу пример выводит
в Serial текст ответа. Адрес в примере \emph{условный}: сервиса по нему может не быть. Замените его на
адрес своего сервиса, у которого есть открытый API; многие метеосервисы выдают JSON бесплатно, но часто
требуют регистрации и ключ.

\codecap{GET-запрос и вывод ответа}
\codeinput{l18_http_get}

Положительный код --- ответ сервера (200 успех, 404 страницы нет, 500 сбой на сервере). Отрицательный
код --- ошибка на самой плате: нет соединения, истёк тайм-аут. Функция \code{HTTPClient::errorToString()}
превратит код в понятное сообщение.

\begin{caution}
Адрес вида \code{http://} передаёт всё открытым текстом. Не отправляйте так пароли и ключи доступа.
Для защищённого канала нужен HTTPS с проверкой сертификата (урок 13). Кроме того, не отправляйте запросы
чаще, чем разрешено сервису: обычно хватает раза в минуту или реже, иначе ключ заблокируют.
\end{caution}

\subsection{JSON и ArduinoJson}

JSON --- текстовый формат для данных: пары «имя: значение» в фигурных скобках. Он понятен и людям, и
программам.
\begin{figure}[H]
\centering
\begin{tikzpicture}[font=\sffamily\scriptsize]
  \node[draw,rounded corners=2pt,fill=blkmem,align=left,font=\ttfamily\small,inner sep=6pt] (j) at (0,0)
    {\{\\ \ \ "city": "Moscow",\\ \ \ "temp": -3.5,\\ \ \ "humidity": 80\\ \}};
  \node[blk,fill=blkcpu,minimum width=4.4cm,anchor=west] (v1) at (7.8,0.9)  {\code{const char* city}};
  \node[blk,fill=blkcpu,minimum width=4.4cm,anchor=west] (v2) at (7.8,0)    {\code{float temp}};
  \node[blk,fill=blkcpu,minimum width=4.4cm,anchor=west] (v3) at (7.8,-0.9)  {\code{int humidity}};
  \draw[arr,accent] (j.east) ++(0,0.65) -- (v1.west);
  \draw[arr,accent] (j.east) ++(0,0.1)  -- (v2.west);
  \draw[arr,accent] (j.east) ++(0,-0.45) -- (v3.west);
  \node[above=2pt,font=\sffamily\small] at (4.9,0.05) {\code{doc["temp"]}};
\end{tikzpicture}
\caption{Разбор JSON: текст ответа превращается в обычные переменные программы}
\end{figure}

Библиотека \textbf{ArduinoJson} (версия 7, ставится через Library Manager) умеет разобрать такой текст.
Создайте объект \code{JsonDocument}, вызовите \code{deserializeJson()} и \emph{обязательно} проверьте результат: сервер может
вернуть страницу с ошибкой вместо JSON. Оператор \code{|} подставляет значение по умолчанию, если поля в ответе нет.

\codecap{Разбор JSON-ответа}
\codeinput{l18_json}

\begin{tip}
Если разбор не удаётся, сначала выведите в Serial сырой ответ \code{http.getString()}. Чаще всего вместо JSON приходит
HTML-страница ошибки или перенаправление.
\end{tip}

\subsection{Точное время по NTP}

В плате нет часов с батарейкой: после включения она «не знает», какое сейчас число. Узнать его можно
у серверов времени по протоколу \textbf{NTP}. Функция \code{configTzTime()} задаёт часовой пояс и серверы,
после чего время синхронизируется само в фоне. Часовой пояс записывается в формате POSIX, где знак
\emph{обратный} привычному: Москва (UTC+3, без перехода на летнее время) записывается как \code{"MSK-3"}.

Функция \code{getLocalTime(\&t, timeout)} ждёт синхронизации не дольше заданного времени и возвращает
\code{false}, если время ещё не получено. Это нужно и для HTTPS (урок 13): без правильной даты проверка
срока сертификата невозможна.

\codecap{Московское время в Serial Monitor}
\codeinput{l18_ntp}

\subsection{LittleFS: файлы во Flash}

В уроке 11 мы сохраняли настройки в NVS, в уроке 10 --- журнал на SD-карту. \textbf{LittleFS} занимает
промежуточное место: это настоящая файловая система внутри Flash платы. Файлы лежат в своём разделе
памяти, их можно открывать, читать и дописывать, как на карте, и никакие дополнительные детали не нужны.

\begin{center}\small
\begin{tabularx}{\textwidth}{@{}l X X X@{}}
\toprule
 & NVS (урок 11) & LittleFS & SD-карта (урок 10) \\
\midrule
Что хранит & пары «ключ--значение» & файлы и папки & файлы и папки \\
Объём & десятки килобайт & около \SIrange{1}{1.5}{MB} (зависит от схемы разделов) & гигабайты \\
Нужны ли детали & нет & нет & модуль и карта \\
Хорош для & небольшие настройки, счётчики & страницы, конфиги, небольшие журналы & большие журналы, фото \\
\bottomrule
\end{tabularx}
\end{center}

\codecap{Чтение и запись файла настроек}
\codeinput{l18_littlefs}

Параметр \code{true} в \code{LittleFS.begin()} форматирует раздел при первом запуске. Так же, как
и на SD, файл нужно закрывать после записи. Помните про износ Flash: не пишите в файл на каждом
проходе \code{loop()}, как и с NVS.

\subsection{Самостоятельный опыт: сетевые часы}

Соберите отдельный скетч: плата подключается к Wi-Fi, получает московское время и печатает его раз в секунду.
Затем добавьте в LittleFS файл \code{/boot.txt}, куда при каждом запуске дописывается строка с датой и временем
старта. Перезагрузите плату несколько раз и прочитайте файл целиком: вы увидите историю запусков.

\begin{trouble}
\begin{itemize}[leftmargin=*]
  \item \textbf{Код ответа отрицательный} --- нет Wi-Fi или интернета, неверный адрес, либо сайт требует HTTPS, а в адресе стоит \code{http://}.
  \item \textbf{\code{deserializeJson()} возвращает ошибку} --- сервер прислал не JSON (страница ошибки, пустой ответ); выведите сырой текст и проверьте адрес и ключ доступа.
  \item \textbf{Время не приходит или показывает 1970 год; \code{LittleFS.begin()} не работает} --- закрыт доступ к NTP или нет интернета; для LittleFS выберите схему разделов с файловой системой.
\end{itemize}
\end{trouble}

\begin{summary}
\begin{itemize}[leftmargin=*]
  \item \code{HTTPClient} отправляет GET-запрос; положительный код --- ответ сервера, отрицательный --- ошибка сети.
  \item JSON разбирают через \code{deserializeJson()}, обязательно проверяя ошибку и подставляя значения по умолчанию.
  \item Время приходит по NTP через \code{configTzTime()} (для Москвы \code{"MSK-3"}), а файлы настроек удобно хранить в LittleFS.
\end{itemize}
\end{summary}

\begin{selfcheck}
\begin{enumerate}
  \item Что происходит между вызовом \code{http.GET()} и получением ответа?
  \item Почему результат \code{deserializeJson()} нужно проверять?
  \item Чем LittleFS отличается от NVS и от SD-карты?
\end{enumerate}
\end{selfcheck}

\begin{tasks}
\begin{enumerate}
\item \lv{1} Выполните GET-запрос и выведите код ответа и размер данных.
\item \lv{1} Разберите готовый JSON и выведите одно из его полей.
\item \lv{2} Выведите московское время на OLED в формате ЧЧ:ММ:СС с обновлением раз в секунду.
\item \lv{2} Показывайте температуру из интернета на OLED вместе со временем получения.
\item \lv{3} Раз в десять минут запрашивайте данные и сохраняйте последнее значение в LittleFS, чтобы показывать его без сети.
\item \lv{4} Сделайте журнал с отметками времени в LittleFS, который не растёт бесконечно: при превышении размера старые записи удаляются.
\end{enumerate}
\end{tasks}
